% !TITLE 饮马长城窟行
% !AUTHOR 无名氏
% !DYNASTY 两汉
% !GENRE 乐府诗
% !ORDER 6

\begin{poem}
青青河畔草，绵绵思远道。
远道不可思，宿昔\textsuperscript{1}梦见之。
梦见在我傍\textsuperscript{2}，忽觉在他乡。
他乡各异县，辗转不相见。
枯桑知天风，海水知天寒。
入门各自媚\textsuperscript{3}，谁肯相为言\textsuperscript{4}？
客从远方来，遗我双鲤鱼\textsuperscript{5}。
呼儿烹鲤鱼，中有尺素书\textsuperscript{6}。
长跪\textsuperscript{7}读素书，书中竟何如？
上言加餐食，下言长相忆。
\end{poem}

\begin{annotations}
\notetext{1}{宿昔：这里指昨夜。宿，昨夜；昔，往日。}
\notetext{2}{傍：旁边、身侧。}
\notetext{3}{各自媚：各自亲近、体贴对方。这里形容别家夫妻在房中团聚恩爱的温馨场景。}
\notetext{4}{相为言：为我说话、安慰我。相，指代“我”；言，安慰、问候。}
\notetext{5}{双鲤鱼：藏信的木匣。因雕刻成两只鲤鱼相对相合的形状而得名，是古人寄托深情和传达书信的器物。}
\notetext{6}{尺素书：写在白绢上的简短家书。尺素，一尺长的白绢（古代多作为书写材料）。}
\notetext{7}{长跪：古代的一种庄重坐姿，双膝着地，直立上身。这里形容妻子极其郑重、恭敬地捧读丈夫来信。}
\end{annotations}

\begin{appreciation}
《饮马长城窟行》是汉乐府中写相思的名篇之一，“鱼传尺素”这个典故就是从这首诗来的。它写一个妻子思念远征的丈夫：从眼前的青草，到梦里的相会，再到一封家书，两千年前的思念，被它写全了。

青青河畔草，绵绵思远道。河边的草长得绵延一片，“绵绵”既说草，也说思念——她看到什么，都往那条路上想。远道不可思，宿昔梦见之。宿昔是昨夜。路太远，想也想不着，只能在梦里见。梦见在我傍，忽觉在他乡。梦里他就在身边，一睁眼，人还在千里之外。这一睁眼的落差，大概是全诗最难受的一刻。还记得《少司命》里的“悲莫悲兮生别离”吗？活着却被生生分开：两个人都在，中间却隔着永远走不完的路。

枯桑知天风，海水知天寒。枯桑还知道风吹，海水还知道天寒——她的心没有麻木，冷暖感知得清清楚楚，正因为清楚，才格外冷。入门各自媚，谁肯相为言？别人家夫妻进门，各自亲亲热热，谁肯来跟我说句话、安慰我一句？别人的团圆，衬出她的冷清。

客从远方来，遗我双鲤鱼。双鲤鱼不是两条鱼，是刻成两条鲤鱼模样的木匣子，里头装信。呼儿烹鲤鱼，中有尺素书。“烹鲤鱼”就是打开这个木匣——古人收到信，要先“杀鱼取信”，这一套很有仪式感。尺素是一尺长的白绢，古代的信纸。

长跪读素书，书中竟何如？长跪是古时最恭敬的坐姿，双膝跪地、挺直上身。读丈夫一封信，郑重成这个样子。上言加餐食，下言长相忆。信的开头说：好好吃饭。结尾说：我一直想着你。没有一句情话，就这两件事。可隔着千山万水的人，最深的牵挂，恰恰就是这两件。

《行行重行行》的最后一句是什么？“努力加餐饭”。这一首的结尾是“上言加餐食，下言长相忆”。两个无名诗人，隔着几十年的光阴，不约而同，把爱藏进了同一件事：好好吃饭。后来杜甫在战乱里写“家书抵万金”，说的也是这种信。今天的人掏出手机就能发消息，很难想象一封要“杀鱼”才取得出的信有多重。

所以妈妈说“好好吃饭”的时候，那不是客套——是她写给你的尺素书。将来你出门在外，记得回她一句：吃了。
\end{appreciation}
